\documentclass[11pt, a4paper]{article}
\usepackage[utf8]{inputenc}
\usepackage[T1]{fontenc}
\usepackage{geometry}
\usepackage{hyperref}
\usepackage{parskip}

\geometry{left=1in, right=1in, top=1in, bottom=1in}

\title{\textbf{A Letter to Gilbert Strang: On Rank-One Matrices, National Flags, and the First Token in My Concept Space}}
\author{Yiwei Zhang (\textsc{Y\v{i}w\`{e}i}, \texttt{math4mad}) \\ \small The Cora Project --- \texttt{github.com/math4mad}}
\date{\today}

\begin{document}

\maketitle

Dear Professor Strang,

I am writing to you not as a student seeking instruction, but as someone whose entire intellectual journey in machine learning and artificial intelligence began with a single anchor point in your work.

Fifteen years ago, I asked a question on Zhihu (a Chinese Q\&A platform): ``If I want to learn linear algebra now, what should I study?'' Someone recommended your MIT OpenCourseWare. For many, this may have been a routine suggestion. For me, it was the first \textbf{token} that activated a vast region of my concept space. I keep the date vague on purpose --- somewhere around 2011 --- because memory, as this letter's framework insists, decays in its date column while its topology of events stays crisp: the timestamp is degradable; what matters is what happened.

I had taken linear algebra in my second year of university and ranked second in my class. But I will be honest: at that time, linear algebra in my mind was ``like shit.'' I could solve problems, but I did not \textit{see} anything. It was only through years of grinding---through SVD, through face recognition, through JPEG compression---that the formulas gradually acquired texture and meaning.

Two concepts from your \textit{Introduction to Linear Algebra} --- whose 5th edition sits on my shelf today, bought for reasons entirely its own --- became the bedrock of my understanding:

\begin{enumerate}
    \item \textbf{Rank-One Matrices.} You taught me that every matrix is a sum of rank-one matrices---that these are the fundamental building blocks, the atoms of linear algebra.
    \item \textbf{National Flags as Rank-One Encodings.} You used the stripes of a national flag to demonstrate how rank-one matrices encode visual information---each colored stripe is a rank-one matrix, and the entire flag is the superposition of these rank-one components. This image---a flag decomposed into its rank-one stripes---was the moment I finally \textit{saw} linear algebra, not just computed it.
\end{enumerate}

These two ideas---rank-one decomposition and the flag example---became the first coordinate system in my cognitive map. They were the anchor, the first token, the probe that activated everything that followed.

Recently, I formalized a theory called the \textbf{Atlas Framework}, which proposes that tokens in large language models are not carriers of information but \textbf{probes into a high-dimensional concept space}. The theory unifies face recognition (eigenfaces and coefficient strength), JPEG encoding (DCT basis and compression), synaptic memory (Hebbian learning and coefficient decay), and semantic inference (token embeddings as coefficient vectors) under a single mathematical structure: \textbf{representation via basis coefficients}.

The key insight is that all four domains are related by \textbf{basis changes}---the same underlying information represented in different coordinate systems. A token provides the coefficients; the model's weights provide the basis; the output is the reconstruction.

This theory emerged from a simple observation: when I showed the number ``11'' to my niece, she saw only a number. When I showed it to my brother, a former basketball fan, he saw Yao Ming. The same token, different concept spaces. The probe did not transmit data; it activated a coordinate in a semantic network whose structure was determined by the receiver's prior experience.

I am sharing this with you because your rank-one matrices and the national flag example were the foundation upon which this entire edifice was built. Without your teaching---without that first token activated in my concept space fifteen years ago---I would never have arrived here.

You have spent your career helping people \textit{see} linear algebra, not just compute it. I am one of those people. And what I have seen, through the lens you provided, is that the same mathematical structure --- basis, coefficient, and the miracle that a few atoms suffice --- governs how we recognize faces, compress images, store memories, and navigate meaning.

Thank you for being the first anchor in my concept space. Fifteen years later, when a child answered ``11 with ``just a number and her uncle answered with ``Yao Ming, it was your rank-one stripes that let me hear those two answers as different bases of one space.

With deep respect,

Yiwei Zhang (\textsc{Y\v{i}w\`{e}i}, \texttt{math4mad}) \\
\small The Cora Project --- \texttt{github.com/math4mad} \quad\texttt{https://math4mad.github.io/cora-atlas/}

\end{document}